% BHU PYQ -- Manually transcribed & error-corrected
% Paper: GLB-201: Elements of Mineralogy and Crystallography
% Exam: B.Sc. (Hons.) Semester II Examination, 2023-24

\documentclass[11pt,a4paper]{article}
\usepackage[a4paper,top=2.5cm,bottom=2.5cm,left=2.8cm,right=2.8cm,headheight=14pt]{geometry}
\usepackage{amsmath,amssymb}
\usepackage[T1]{fontenc}
\usepackage[utf8]{inputenc}
\usepackage{lmodern,microtype,fancyhdr,enumitem,hyperref,lastpage,parskip}

\hypersetup{colorlinks=true,urlcolor=black,linkcolor=black,
  pdftitle={GLB-201 Elements of Mineralogy and Crystallography -- BHU B.Sc. Sem II 2023-24}}
\pagestyle{fancy}\fancyhf{}
\renewcommand{\headrulewidth}{0.4pt}\renewcommand{\footrulewidth}{0.4pt}
\fancyhead[L]{\small   \textit{Banaras Hindu University}}
\fancyhead[R]{\small   \textit{GLB-201: Elements of Mineralogy \& Crystallography}}
\fancyfoot[C]{\small Page  \thepage\ of \pageref{LastPage}}


\newlist{parts}{enumerate}{2}
\setlist[parts,1]{label=(\alph*),leftmargin=*,itemsep=5pt,topsep=3pt}

\newcommand{\pts}[1]{\hfill   \textit{[#1]}}


\begin{document}

\begin{center}
  {\large\bfseries BANARAS HINDU UNIVERSITY}\\[2pt]
  {
\normalsize B.Sc. (Hons.) --- Semester II Examination, 2023-24}\\[6pt]
  \rule{0.8\linewidth}{0.5pt}\\[6pt]
  {\large\bfseries Paper: GLB-201 --- Elements of Mineralogy and Crystallography}\\[4pt]
  {
\normalsize Time: 3 Hours \hfill Maximum Marks: 70}\\[2pt]
  {\small Ref. No.: 074430}
\end{center}
\rule{\linewidth}{0.4pt}
\smallskip



\noindent   \textit{Instructions: Attempt   \textbf{all} three Sections A, B, and C according to the directions.}
\bigskip

\bigskip


\noindent {\large\bfseries SECTION --- A}
\rule{\linewidth}{0.4pt}\smallskip



\noindent   \textit{Answer   \textbf{any two} questions. (14 marks each)}
\medskip




\noindent   \textbf{Question 1.} Why is the megascopic property of a mineral crucial for its nomenclature? A student described a mineral with the physical property ``Perfect, 2 sets, $\{001\}$''. Which mineral property is he referring to? Explain in detail. \pts{14}

\medskip


\noindent   \textbf{Question 2.} Define the optical properties of minerals and illustrate with diagrams the following properties (   \textbf{any four}): \pts{14}

\begin{parts}
  \item Refractive index
  \item Becke line test
  \item Interference colour
  \item Optical relief
  \item Extinction angle
  \item Pleochroism
\end{parts}

\medskip


\noindent   \textbf{Question 3.} What is the basis for classifying silicate minerals? Discuss the different silicate groups and describe two processes through which quartz and feldspar can form. \pts{14}

\medskip


\noindent   \textbf{Question 4.} Two olivine crystals (a common gemstone: ``Peridot'') were excavated from basaltic rocks at two different depths within an igneous complex. Do they possess the same interfacial angles? Justify your answer with suitable reasons. Under which conditions is a reflecting goniometer used instead of a contact goniometer? Also explain the elements of symmetry in a three-dimensional coordinate reference frame. \pts{14}

\bigskip


\noindent {\large\bfseries SECTION --- B}
\rule{\linewidth}{0.4pt}\smallskip



\noindent   \textit{Answer   \textbf{any two} questions. (14 marks each)}
\medskip




\noindent   \textbf{Question 5.} Write short notes on the following: \pts{$3.5  \times 4 = 14$}

\begin{parts}
  \item What is the Nicol prism and why is it important in a polarizing microscope?
  \item How to check the position of a polarizing microscope ready to examine the optical properties of a mineral.
  \item What are isotropic and anisotropic minerals? How to identify these minerals under a microscope.
  \item Define the optical properties that can distinguish quartz, plagioclase, microcline, and orthoclase with diagrams.
\end{parts}

\medskip


\noindent   \textbf{Question 6.} Briefly explain the following sentences (   \textbf{any four}) with suitable examples wherever necessary: \pts{$3.5  \times 4 = 14$}

\begin{parts}
  \item Pyrite generally breaks into a cubic crystal whereas calcite breaks into a rhombic crystal.
  \item For some minerals, a streak plate reveals their streak, while for some, it indicates its hardness.
  \item Quartz exhibits a range of colours, while plagioclase does not.
  \item Anion or anionic group based classification of minerals is considered the most valid scheme.
  \item Diamond and graphite, despite having the same chemical composition, have different specific gravity.
\end{parts}

\medskip


\noindent   \textbf{Question 7.} Explain the following: \pts{$7  \times 2 = 14$}

\begin{parts}
  \item Why are 5-fold, 7-fold, and 8-fold rotational symmetries not possible in crystals?
  \item In a cubic crystal, is the face $(111)$ a general form or a special form?
\end{parts}

\bigskip


\noindent {\large\bfseries SECTION --- C}
\rule{\linewidth}{0.4pt}\smallskip



\noindent   \textit{All questions are compulsory. (1 mark each)}
\medskip




\noindent   \textbf{Question 8.} Answer in one/two words or sentence(s): \pts{$1  \times 14 = 14$}

\begin{parts}
  \item Write two optical properties that depend on the refractive index.
  \item How to differentiate between cleavage and fracture under a microscope.
  \item Write two essential mineral properties required to measure the extinction angle.
  \item Define optical twinkling with an example of a twinkling mineral.
  \item Write two optical properties to distinguish muscovite and biotite.
  \item Write two different optical properties of hornblende and augite.
  \item Which mineral is named after a scientist?
  \item Who discovered the polarizing microscope, and in which year?
  \item A mineral scratches topaz and shows good parting. Name the mineral.
  \item Write one mineral formed through a metamorphic process showing a bladed habit.
  \item Which instrument is used to distinguish pitchblende and hornblende?
  \item Tiger's eye and cat's eye describe which property of a mineral?
  \item What is a sphenoid?
  \item What do you mean by a solid angle?
\end{parts}

\vfill\rule{\linewidth}{0.4pt}

\begin{center}\small
\href{https://bhu-exam.vercel.app/}{Click here to visit our website}\\[4pt]
\href{https://me-aryan.vercel.app/}{Click here to visit our contributor}\\[4pt]
\href{https://quashan-me.vercel.app/}{Click here to visit our contributor}
\end{center}
\end{document}
